\documentclass[11pt]{article}
\usepackage[T1]{fontenc}
\usepackage{amsmath,amssymb,amsthm,geometry}
\geometry{margin=.9in}\setlength{\emergencystretch}{2em}
\allowdisplaybreaks[2]
\newtheorem{theorem}{Theorem}
\newcommand{\E}{\mathbb E}\newcommand{\dd}{\,d}
\title{Dissipation versus cross-output exchange on the critical path}
\author{E5 continuation: analytic partial result}
\date{4 November 2026}
\begin{document}\maketitle

\paragraph{Result.}
The residual barrier through half response remains unproved.
The strong-response component of the dissipation can be completed
into a nonnegative square, with a vanishing accumulated error.
The remaining comparison is a particular signed correlation between
the weak-target gradient and the transverse-error correction, after
both have been projected away from the strong-response gradient.
An attempt to propagate the needed covariance bound encounters an
explicit signed weak-gate-motion term. No actual exit, and no
violation of the comparison on the critical path, is asserted.

\section{Reached interval and definitions}
Use REC's exact finite-noise notation and the same critical family
$S_0=cq^{1/\omega}$, with fixed $c>0$,
$\lambda=\rho^2$, $\rho\in[13/20,7/10]$,
$\omega=(1-\lambda)/2$. The model is
\[
 f(D)=\int W(U^TD)_+\dd l,\quad |U|=|W|,\quad
 P_1=N(e_1,q^2I),\quad P_2=N(\rho e_2,q^2I).
\]
The mixture weights are $1/2$ and the loss is
$\mathcal L=\frac14\sum_p\E_{P_p}|D-f(D)|^2$.
Write $S=\int|U|^2\dd l$, $Q_1=1+q^2$, $Q=\lambda+q^2$,
and define the four actual regressions
\[
 m=\E_{P_1}D_1f_1/Q_1,\quad z=\E_{P_1}D_1f_2/Q_1,\quad
 h=\E_{P_2}D_2f_1/Q,\quad n=\E_{P_2}D_2f_2/Q.
\]
Let $g_m,g_z,g_h,g$ be their full two-layer parameter gradients,
in that order, and set $w=Q(1-n)$.
Then $\dot\Theta=(w/2)g+N$, where
\[
 \mathcal R=4\mathcal L-Q(1-n)^2,\qquad N=-\nabla\mathcal R/4,
 \qquad \mathcal R'=-4\|N\|^2-2w\langle g,N\rangle.
\]
We work on the \emph{reached stopped interval} $[\tau_0,\tau]$
of TBW/RQT: $n(\tau_0)=\eta_0$, $\mathcal R(\tau_0)\le r_*^2/16$,
$r_*=\lambda/800$, and $\tau$ is the first subsequent hit of
$n=1/2$ or $\mathcal R=r_*^2$. The proved statements are
\begin{equation}\label{an04pcdv1:interval}
 \tau-t_h\le C_c(1+\ell),\qquad S\le C_c(1+\ell),
 \qquad \ell=\log(1/q).
\end{equation}
No half-response arrival is presumed. On this interval
$m\ge1-r_*/\sqrt{Q_1}>1/2$, $w\ge Q/2>0$,
and all four residual coefficients are bounded by loss decrease.

REC's exact shape decomposition reads
\begin{equation}\label{an04pcdv1:split}
 2N=Q_1(1-m)g_m-Q_1z g_z-Qh g_h+2N_{\rm sh},
 \qquad \|N_{\rm sh}\|\le Cq\sqrt S(1+S).
\end{equation}
It follows from the orthogonal within-cluster residuals, whose
total squared norm is at most $Cq^2(1+S)^2$.
This estimate uses the network's global Lipschitz constant $S$,
and is valid without an angular chart.

\section{Removing the strong-response direction without losing dissipation}
Let $A=\|g_m\|^2$, $e=g_m/\sqrt A$, and let
$P=I-e\otimes e$ denote orthogonal projection in parameter space.
Positive two-homogeneity gives
$\langle g_m,\Theta\rangle=2m$ and $\|\Theta\|^2=2S$.
Consequently
\begin{equation}\label{an04pcdv1:nondegenerate}
 A\ge2m^2/S>0.
\end{equation}
Thus this projection is defined on the entire stopped interval;
there is no inverse of a possibly zero cross-gradient matrix here.
Define
\[
 d=\langle g,e\rangle,\quad \alpha=2\langle N,e\rangle,
 \quad V=-2PN,\quad \beta=\alpha+wd/2.
\]
Orthogonality alone proves the exact identity
\begin{equation}\label{an04pcdv1:exact}
 \mathcal R'=w\langle Pg,V\rangle-\|V\|^2-\beta^2
                         +w^2d^2/4.
\end{equation}
Indeed $2N=\alpha e-V$, so
$4\|N\|^2=\alpha^2+\|V\|^2$ and
$2\langle g,N\rangle=\alpha d-\langle Pg,V\rangle$.
Complete the square in $\alpha$.

REC proves $|\langle g,g_m\rangle|\le CqS$.
Together with \eqref{an04pcdv1:nondegenerate}, this gives
\begin{equation}\label{an04pcdv1:smallparallel}
 |d|\le CqS^{3/2},\qquad
 \int_{\tau_0}^{\tau}w^2d^2\dd t\le C_cq^2(1+\ell)^4=o(1).
\end{equation}
The strong-response adjustment has therefore been accounted for
as a nonnegative square and a paid error. Its mixed dissipation
terms were not discarded.

To expose the two cross regressions, put
\[
 p=\binom{Q_1z}{Qh},\qquad
 H_{ij}=\langle Pg_i,Pg_j\rangle,\qquad
 k_i=\langle Pg,Pg_i\rangle\quad(i,j\in\{z,h\}),
\]
and define
$V_0=Q_1zPg_z+QhPg_h$,
$\mathcal D_\perp=p^THp=\|V_0\|^2$,
$\mathcal C_\perp=p^Tk=\langle Pg,V_0\rangle$.
The correlation $\mathcal C_\perp$ is signed.

\begin{theorem}[Exact remaining comparison on the initialized stopped path]
\label{an04pcdv1:main}
On $[\tau_0,\tau]$,
\begin{equation}\label{an04pcdv1:remaining}
 \boxed{\quad\mathcal R'=w\mathcal C_\perp
                 -\mathcal D_\perp-\beta^2+\mathcal E_q.\quad}
\end{equation}
The error is explicitly, with $E=PN_{\rm sh}$,
\[
 \mathcal E_q=-2w\langle Pg,E\rangle+4\langle V_0,E\rangle
                  -4\|E\|^2+w^2d^2/4,
\]
and obeys
\begin{equation}\label{an04pcdv1:error}
 \int_{\tau_0}^{\tau}|\mathcal E_q|\dd t
 \le C_c\{q(1+\ell)^3+q^2(1+\ell)^4\}=o(1).
\end{equation}
Furthermore, for REC's moments $C_{12},T_2$,
\begin{equation}\label{an04pcdv1:source}
 \mathcal C_\perp=Q_1zC_{12}+QhT_2+O(qS(1+S)).
\end{equation}
The exact upper-face condition is
\begin{equation}\label{an04pcdv1:face}
 w\mathcal C_\perp\le\mathcal D_\perp+\beta^2-\mathcal E_q
       \qquad(\mathcal R=r_*^2,\ n<1/2).
\end{equation}
No assertion that this inequality holds is part of the theorem.
\end{theorem}
\begin{proof}
Equation \eqref{an04pcdv1:split} implies $V=V_0-2E$.
Substitution in \eqref{an04pcdv1:exact} proves
\eqref{an04pcdv1:remaining} with the stated signed error.
Use $\|Pg\|,\|V_0\|\le C\sqrt S$ and
$\|E\|\le Cq\sqrt S(1+S)$ to bound the error pointwise by
$C[qS(1+S)+q^2S(1+S)^2+q^2S^3]$.
Integration with \eqref{an04pcdv1:interval} proves
\eqref{an04pcdv1:error}.
For each cross gradient,
\[
 k_i=\langle g,g_i\rangle-d\langle e,g_i\rangle,
 \qquad |d\langle e,g_i\rangle|\le CqS^2.
\]
REC gives $\langle g,g_z\rangle=C_{12}+O(qS)$ and
$\langle g,g_h\rangle=T_2$. This proves
\eqref{an04pcdv1:source}; the face condition follows from the exact
derivative. No pointwise sign was assigned to an error bound.
\end{proof}

If the residual event occurs first, integration gives the necessary
payment
\begin{equation}\label{an04pcdv1:exit}
 \int_{\tau_0}^{\tau}(w\mathcal C_\perp-\mathcal D_\perp)\dd t
 \ge15r_*^2/16+\int_{\tau_0}^{\tau}\beta^2\dd t-o(1).
\end{equation}
This is conditional on an actual exit; it does not establish one.
The unknown has been localized to the weak-target correlation with
the projected transverse correction, relative to its own dissipation.

\section{First attempted closure: joint Gram control}
The projection matters even though the unprojected $g_z,g_h$
are nearly orthogonal. To see the mixed term explicitly, let
\[
 T_1=\langle g_m,g_z\rangle,
 \quad B_m=\langle g_m,g_h\rangle,
 \quad G_{ij}=\langle g_i,g_j\rangle.
\]
Then
\begin{equation}\label{an04pcdv1:schur}
 \mathcal D_\perp=
 p_z^2G_{zz}+2p_zp_hG_{zh}+p_h^2G_{hh}
                 -\frac{(p_zT_1+p_hB_m)^2}{A}.
\end{equation}
Here $G_{zh}=O(qS)$, $B_m=C_{12}+O(qS)$, and
$T_1=\int W_1W_2|v_1(U)|^2\dd l$ in REC's Gaussian features.
Thus subtracting the strong-response direction produces the
cross-correlation $-T_1B_m/A$ in $H_{zh}$. Replacing
\eqref{an04pcdv1:schur} by its unprojected diagonal terms would
overestimate the available dissipation.

The sharp Cauchy--Schwarz estimate using the two cross gradients is
\begin{equation}\label{an04pcdv1:cs}
 |\mathcal C_\perp|\le\sqrt{\mathcal D_\perp\,\kappa},
 \qquad \kappa=k^TH^\dagger k\le\|Pg\|^2.
\end{equation}
This remains valid at singular $H$: a null vector of $H$ represents
the zero parameter vector and is orthogonal to $k$, so $k$ is in
the range of $H$. Equivalently, $\kappa$ is the squared norm of
the projection of $Pg$ onto the span of $Pg_z,Pg_h$.
It is used pointwise; no inverse is differentiated or integrated.
Consequently
\[
 w\mathcal C_\perp-\mathcal D_\perp
 \le w\sqrt{\mathcal D_\perp\kappa}-\mathcal D_\perp
 \le w^2\kappa/4.
\]
This can be positive. For example, the sufficient comparison
$\mathcal D_\perp\ge w^2\kappa$ would pay the two displayed
terms, but is not supplied by the residual cap or by the reached
small-hit estimates. The actual signed direction could also help;
that direction is exactly the unpaid $\mathcal C_\perp$ in
\eqref{an04pcdv1:face}. This algebra is not a counterexample
on the critical trajectory.

\section{Second attempted closure: transporting the covariances}
We compute the terms that a moment argument would have to control.
For $Q_2=Q$, define the exact Gaussian features
\[
 v_p(U)=Q_p^{-1}\E_{P_p}[D_pD\mathbf1_{U^TD>0}],\qquad
 \theta_p(U)=Q_p^{-1}\E_{P_p}[D_p^2\mathbf1_{U^TD>0}].
\]
They satisfy $v_p=\theta_pe_p+O(q)$ and $0\le\theta_p\le1$
uniformly over every input direction.
The residual projection used in REC gives the pointwise receiver
matrix, with both mixture halves retained,
\[
 R_U=\tfrac{Q_1}{2}[(1-m)e_1-ze_2]v_1^T
       +\tfrac Q2[-he_1+(1-n)e_2]v_2^T+O(q(1+S)).
\]
The error follows by Cauchy--Schwarz on the shape residuals;
it is uniform over receiver gates.

For $C_{12}=\int(U_1)_+(U_2)_+\dd l$, the a.e. chain rule gives
\begin{align}
 2C_{12}'={}&Q_1\int_{++}\theta_1U_2[(1-m)W_1-zW_2]\dd l\nonumber\\
 &+Q\int_{++}\theta_2U_1[-hW_1+(1-n)W_2]\dd l
                  +O(qS(1+S)),\label{an04pcdv1:overlap}
\end{align}
where $++=\{U_1>0,U_2>0\}$ at the current time.
The positive-part product vanishes at each face, so this identity
does not drop a boundary-flux term. At intersections of the faces
the same chain rule applies; no invariant quadrant is assumed.
In particular the target regenerates overlap through
$w\int_{++}\theta_2U_1W_2\dd l$.
That restricted covariance is not the measured total $z$.
The cap bounds the total regression, and does not exclude signed
cancellation with labels outside this overlap region.

There is an additional issue for
$T_2=\int W_1W_2|v_2(U)|^2\dd l$.
Its exact derivative is
\begin{equation}\label{an04pcdv1:transport}
 T_2'=\int(W_1'W_2+W_1W_2')|v_2|^2\dd l+\mathfrak G_2,
 \qquad
 \mathfrak G_2=2\int W_1W_2\,v_2^TDv_2(U)[U']\dd l.
\end{equation}
Thus a moving Gaussian gate contributes even after the ordinary
output derivatives have been included.

Its coefficient can be calculated exactly, without numerics.
For $a_U=U/|U|$ let $t_U=(-a_{U,2},a_{U,1})$. For $U\ne0$,
\begin{equation}\label{an04pcdv1:hessian}
 Dv_2(U)=\gamma_q(U)t_Ut_U^T,\qquad
 \gamma_q(U)=\frac{\rho\,t_{U,2}^2(\lambda t_{U,2}^2+3q^2)}
                         {Qq|U|}
                 \phi(\rho a_{U,2}/q)\ge0.
\end{equation}
Indeed differentiate the Gaussian integral for $v_2$. The surface
$a_U^TD=0$ has density $\phi(\rho a_{U,2}/q)/(q|U|)$ in the
variable $U^TD$. On that surface $D=t_UT$, with
$T\sim N(\rho t_{U,2},q^2)$. Its contribution is
$t_{U,2}\E T^3\,t_Ut_U^T$ and
$t_{U,2}\E T^3=\rho t_{U,2}^2(\lambda t_{U,2}^2+3q^2)$.
This proves the formula, including directions on the weak gate.
Positive initialized radii cannot vanish at a finite time because
the exact flow has uniformly bounded receiver matrices by loss
decrease. Therefore no zero-input label is omitted.

The remaining gate-motion correlation is explicitly
\begin{equation}\label{an04pcdv1:gate}
 \boxed{\quad\mathfrak G_2=
 2\int\gamma_q(U)W_1W_2
           (v_2^Tt_U)(U'^Tt_U)\dd l.\quad}
\end{equation}
Positivity of $\gamma_q$ does not sign the other three factors.
Using only $\|R_U\|\le C$ and balance gives
\[
 |\mathfrak G_2|\le\frac Cq\int|W_1W_2|
                 \phi(\rho a_{U,2}/q)\dd l\le CS/q.
\]
This upper estimate is far too large to close the residual budget
on a logarithmic interval. No sharper signed or integrated bound
on this gate-layer covariance has been proved here for the actual
critical trajectory. Initial alignment and fixed-resident-time
physical bounds do not provide it at the moving small hit or
through the subsequent stopping interval.

\paragraph{Stop and quantifiers.}
The joint Gram attempt leaves \eqref{an04pcdv1:face} unpaid;
the covariance-transport attempt leaves the restricted overlap
covariance in \eqref{an04pcdv1:overlap} and the signed gate-motion
term \eqref{an04pcdv1:gate} uncontrolled. We stop after these two
approaches. Neither original source $zC_{12}$ or $hT_2$ has been
proved independently dominated, and we do not identify one as
the actual cause of an exit that has not been established.
The precise aggregate missing correlation is
$\mathcal C_\perp=\langle Pg,Q_1zPg_z+QhPg_h\rangle$.

The residual face has strict initialized entry, but its outward
sign is unresolved. The half-response boundary is a stopping
target, and a simultaneous event already reaches that response.
No monotonicity or unproved lower response barrier was used.
All constants may depend on fixed $c,\rho$ and the prior entry
choices; no uniformity in $c\to0$ or $c\to\infty$ is claimed.
The new physical estimates apply on the stopped logarithmic
interval by exact flow identities and the proved mass and time
bounds. No COD constant is claimed at $s=t_\eta-t_h$.
This work does not prove half response or a learned probe minimum
below the coherent comparator. The probe's remaining signed
motion would still need control even after this barrier closes.
No previous increment, foundation, checkpoint, or remote file is
modified. There is no empirical or computer-assisted proof premise.
\end{document}
